% ==============================================================================
% SECCIÓN 9: PLAN DE GESTIÓN DE CAMBIOS
% ==============================================================================
\section{Plan de Gestión de Cambios}

Ante cualquier cambio solicitado, se revisarán los requisitos afectados y se evaluará su impacto en la planificación, el diseño y el desarrollo de la aplicación. Una vez aprobado el cambio, se realizarán las modificaciones necesarias tanto en la aplicación como en la documentación y memoria del proyecto, manteniendo siempre actualizados los requisitos.

\subsection{Procedimiento de Evaluación y Aprobación}
\begin{enumerate}[leftmargin=*]
  \item \textbf{Registro de Solicitud:} Registro formal de la propuesta de modificación identificando origen y justificación.
  \item \textbf{Análisis de Impacto:} Los coordinadores y el Gestor de Calidad evaluarán la repercusión sobre el cronograma en ClickUp, el esfuerzo en horas de los miembros, los costes y la accesibilidad.
  \item \textbf{Decisión Formal:} Si el cambio es aprobado por el Coordinador, el Gestor de Calidad y el Moderador, se traslada al Product Backlog y se actualizan los diagramas UML y memorias.
  \item \textbf{Preaviso de 3 Días:} Si la alteración afecta a las fechas de entrega pactadas en el pliego técnico, se comunicará formalmente al profesorado con un mínimo de \textbf{3 días de antelación}.
\end{enumerate}
